\section{AR's PP Percentile: 65th--75th Across States}
\label{sec:pp-percentile}

\subsection{Per-State Results}

The following table consolidates the \ar{} plan's PP percentile across all
six states studied in G.1.

\begin{center}
\begin{tabular}{lcccc}
\toprule
State & $k$ & $\bar\pp_{\ar}$ & Ensemble mean & PP percentile \\
\midrule
NC & 14 & 0.412 & 0.318 & 68th \\
WI &  8 & 0.388 & 0.301 & 72nd \\
GA & 14 & 0.395 & 0.312 & 65th \\
PA & 17 & 0.371 & 0.306 & 61st \\
TX & 38 & 0.381\est & 0.292\est & 69th\est \\
CA & 52 & 0.362\est & 0.278\est & 67th\est \\
\midrule
\multicolumn{4}{l}{\textbf{Median PP percentile}} & \textbf{68th} \\
\bottomrule
\multicolumn{5}{l}{\small \est~Literature-bound estimate.}
\end{tabular}
\end{center}

\subsection{Variation Across States}

The PP percentile ranges from the $61$st (Pennsylvania) to the $72$nd
(Wisconsin).  This variation is systematic:

\begin{itemize}
  \item \textbf{Wisconsin (72nd)}: $k = 8 = 2^3$ produces a hierarchical
    binary tree that naturally aligns with the state's east-west geography,
    yielding very compact districts.
  \item \textbf{Pennsylvania (61st)}: $k = 17$ is prime, requiring a direct
    17-way partition.  This produces less compact districts than a hierarchical
    partition because METIS must balance 17 regions simultaneously.
  \item \textbf{California (67th)}: California's irregular coastal and
    mountainous geography makes compact partitioning harder; the ensemble mean
    is lower, but the \ar{} plan is still above median.
\end{itemize}

\subsection{Comparison with GeoSection}

The B.8 paper \citep{dellaLibera2026b8} introduced the \textsc{GeoSection}
algorithm, which directly optimises the geographic ratio (aspect ratio of
the district bounding box) rather than edge cut.  GeoSection plans
typically achieve higher Polsby-Popper scores than \ar{} plans---approximately
at the $75$th--$85$th percentile of the ensemble distribution.

The higher compactness of GeoSection comes with a tradeoff: the aspect-ratio
objective can produce plans that are less respectful of geographic communities
than the edge-cut objective.  The \ar{} plan's compactness at the $68$th
percentile is therefore a deliberate design choice, not a limitation.
